% Good stuff ... put in cover letter

Declining sensitivity and synchrony are almost always interpreted as evidence that the long-standing and widespread thermal-sum model is now wrong or incomplete.

In this way, our results reveal the key mechanism that determines whether synchrony declines: the temperature regime under which accumulation occurs. When average temperatures increase day over day, as they often do near the equinox, springtime events tend to be tightly synchronized. But when warming shifts those events closer to the solstice, where average temperatures are comparatively stationary, event times become more variable. The shift in regimes is observable as a `scattered spring' in which, as springtime events advance, they become less synchronized. 

 First, it challenges an extensive body of work that has attempted to revise the thermal-sum model by adding additional factors or mechanisms \citep{fu2015declining,kovaleskipreprint}. These factors can influence the timing of biological events, but our work suggests they do not themselves explain variation in sensitivity and synchrony. Further, they may lead to inaccurate forecasts if they do not themselves properly model the basic thermal sum model. 

To understand the role played by any additional factors, the basic relationship between event time and temperature must be modeled correctly. Indeed, the literature often assumes a linear relationship between temperature and event time with constant variance, which we have shown is inconsistent with the thermal-sum model. Our framework provides a simple and transparent way to encode the thermal-sum model, on top of which complexity may be added and better studied.

The second consequence is that the thermal-sum model predicts a `scattered spring' or general desynchronization of springtime events as their timing moves closer to the winter solstice.